\documentclass[
    a4paper,
    man,
    floatsintext
]{glossaPX2}

% Encoding and Language
\usepackage[T1]{fontenc}
\usepackage[american]{babel}

% Bibliography — biblatex with APA 7 (required by Glossa Psycholinguistics)
\usepackage[
  style=apa,
  backend=biber,
  sorting=nyt
]{biblatex}
\addbibresource{example.bib}

% Tables — needed for pandoc output
\usepackage{longtable}
\usepackage{array}    % pandoc uses >{\raggedright\arraybackslash}p{} column types
\usepackage{calc}     % pandoc uses \real{} expressions for column widths
\usepackage{multirow}

% Other styling
% Note: xcolor must NOT be reloaded (option clash with glossaPX2.cls)
\usepackage[font={footnotesize,it}]{caption}
\usepackage{csquotes}  % required by biblatex APA
\usepackage{booktabs}
\usepackage{tabularx}
\usepackage{graphicx}
\usepackage{cgloss}    % interlinear glosses (\gll, \glt)
\usepackage{amsmath}
\usepackage{ragged2e}
\usepackage{hyperref}
\usepackage{enumitem}
\newcommand{\quotes}[1]{``#1''}
\usepackage[title]{appendix}
\usepackage[section]{placeins}

% Example numbering: load either linguex OR gb4e in your document via header-includes.
% Do NOT load both — they conflict. If using linguex, also add:
%   \renewcommand{\firstrefdash}{}  (removes dash in cross-references)

\usepackage{float}
\usepackage{gb4e}
\makeatletter
\@ifpackageloaded{caption}{}{\usepackage{caption}}
\AtBeginDocument{%
\ifdefined\contentsname
  \renewcommand*\contentsname{Table of contents}
\else
  \newcommand\contentsname{Table of contents}
\fi
\ifdefined\listfigurename
  \renewcommand*\listfigurename{List of Figures}
\else
  \newcommand\listfigurename{List of Figures}
\fi
\ifdefined\listtablename
  \renewcommand*\listtablename{List of Tables}
\else
  \newcommand\listtablename{List of Tables}
\fi
\ifdefined\figurename
  \renewcommand*\figurename{Figure}
\else
  \newcommand\figurename{Figure}
\fi
\ifdefined\tablename
  \renewcommand*\tablename{Table}
\else
  \newcommand\tablename{Table}
\fi
}
\@ifpackageloaded{float}{}{\usepackage{float}}
\floatstyle{ruled}
\@ifundefined{c@chapter}{\newfloat{codelisting}{h}{lop}}{\newfloat{codelisting}{h}{lop}[chapter]}
\floatname{codelisting}{Listing}
\newcommand*\listoflistings{\listof{codelisting}{List of Listings}}
\makeatother
\makeatletter
\makeatother
\makeatletter
\@ifpackageloaded{caption}{}{\usepackage{caption}}
\@ifpackageloaded{subcaption}{}{\usepackage{subcaption}}
\makeatother

% pandoc: scale images to fit line width, preserving aspect ratio
\makeatletter
\newsavebox\pandoc@box
\newcommand*\pandocbounded[1]{%
  \sbox\pandoc@box{#1}%
  \Gscale@div\@tempa{\linewidth}{\wd\pandoc@box}%
  \ifdim\@tempa\p@<\p@\scalebox{\@tempa}{\usebox\pandoc@box}\else#1\fi%
}
\makeatother

% pandoc: tight lists
\providecommand{\tightlist}{%
  \setlength{\itemsep}{0pt}\setlength{\parskip}{0pt}}

\title[Instructions for authors]{Instructions for Glossa
Psycholinguistics Submissions}

  \author[Smith \& Jones]{
    \spauthor{
  Alice Smith\\
            \institute{University of Nowhere}\\
            \small{alice.smith@nowhere.edu\\}
        \small{ORCID: 0000-0000-0000-0000}
    }
   \AND
    \spauthor{
  Bob Jones\\
            \institute{Institute of Made-Up Science}\\
            \small{bob.jones@madeup.edu\\}
      }
    }

\begin{document}
\maketitle


\begin{abstract}
This document contains instructions for preparing a submission to Glossa
Psycholinguistics using the \textbf{glossapsycho} Quarto template. It
also serves as a demonstration of the template's formatting: the
document you are reading is itself typeset using the template. Replace
the YAML fields above and the section content below with your own
manuscript.
\end{abstract}

\begin{keywords}
template; reference; example
\end{keywords}

\section{Installation}\label{installation}

The easiest way to install this template is directly from GitHub. Run
the following command from your project directory:

\begin{verbatim}
quarto add znhoughton/quarto-templates/glossapsycho
\end{verbatim}

Alternatively, if you have a local clone of the repository, you can
install from the local path instead:

\begin{verbatim}
quarto add "path/to/quarto-templates/glossapsycho"
\end{verbatim}

Either method copies the \texttt{\_extensions/glossapsycho/} folder into
your project. Once installed, set the format in your document's YAML
front matter:

\begin{verbatim}
format:
  glossapsycho-pdf: default
\end{verbatim}

Then render with:

\begin{verbatim}
quarto render your-document.qmd
\end{verbatim}

The extension folder is self-contained and can be committed to your
repository so that collaborators do not need a separate installation
step. The \texttt{\_extensions/} folder should not be listed in
\texttt{.gitignore} if you want the template to be available to everyone
who clones the repository. To update the template to the latest version
at any time, simply re-run the \texttt{quarto\ add} command from above.

\section{Text Formatting}\label{text-formatting}

This document is typeset using the \texttt{glossapsycho} template. The
body text demonstrates the default font and spacing defined by the
bundled \texttt{glossaPX2.cls} class file. Inline formatting options
include \textbf{bold}, \emph{italic}, \texttt{inline\ code}, and
\textsc{small caps}.

Paragraphs are separated by a blank line. A single line break within a
paragraph does not produce a new paragraph. Line spacing and page
geometry are fixed by the class file and should not be overridden. Do
not add custom geometry, font, or spacing commands to your document.

Section headings use \texttt{\#} for top-level sections and
\texttt{\#\#} for subsections. The template numbers sections
automatically; do not add manual numbering to heading text
\autocite{Smith2020}.

\section{Tables}\label{tables}

Use an R chunk with \texttt{knitr::kable()} and the
\texttt{\#\textbar{}\ label:} and \texttt{\#\textbar{}\ tbl-cap:} chunk
options to produce tables.

\begin{table}

\centering{

\begin{tabular}{lrr}
\toprule
Condition & Mean & SD\\
\midrule
A & 1.2 & 0.3\\
B & 2.4 & 0.5\\
\bottomrule
\end{tabular}

}

\caption{\label{tbl-example}A simple table showing condition means.}

\end{table}%

Cross-reference with \texttt{@tbl-example}, which produces
Table~\ref{tbl-example}. The \texttt{booktabs} and \texttt{tabularx}
packages are loaded by the template and available for more complex table
formatting. Keep tables narrow enough to fit within the text width; the
single-column layout accommodates wider tables than two-column formats.

\section{Figures}\label{figures}

Use an R chunk with \texttt{\#\textbar{}\ label:\ fig-xxx} and
\texttt{\#\textbar{}\ fig-cap:} chunk options to produce figures with
captions and cross-reference labels.

\begin{figure}[ht]
  \centering
  \fbox{\rule{0pt}{3cm}\rule{0.8\linewidth}{0pt}}
  \caption{A simple figure placeholder.}
  \label{fig-example}
\end{figure}

Cross-reference with \texttt{\textbackslash{}ref\{fig-example\}}. Figure
captions are set in small italic type by the template's caption style.
Figures are floating by default and will be placed at the top or bottom
of the nearest page. For publication-quality plots, prefer vector
formats (PDF or SVG) over raster formats (PNG, JPEG) to ensure sharpness
at any print resolution.

\section{Equations}\label{equations}

Inline math: \(y = \beta_0 + \beta_1 x + \varepsilon\).

Display equation:

\begin{equation}\phantomsection\label{eq-model}{
Y_{ij} = \beta_0 + \beta_1 X_{ij} + u_i + \varepsilon_{ij}
}\end{equation}

Cross-reference with \texttt{@eq-model}, which produces
Equation~\ref{eq-model}.

Multi-line display equations:

\[
\mu_i = \alpha + \beta x_i, \quad y_i \sim \mathcal{N}(\mu_i, \sigma^2)
\]

Use display equations for any expression that is important enough to be
referred to by number. Number only equations that are actually
cross-referenced in the text; unnumbered display equations use the plain
\texttt{\$\$...\$\$} syntax without a label.

The \texttt{amsmath} package is loaded by the template and provides
environments such as \texttt{align}, \texttt{gather}, and \texttt{cases}
via raw LaTeX blocks. Inline math uses standard \texttt{\$...\$}
delimiters. Avoid mixing LaTeX math commands with Pandoc markdown
constructs in the same expression.

\section{Interlinear Glosses}\label{interlinear-glosses}

Glosses are typeset with \texttt{gb4e} via raw LaTeX blocks. Add
\texttt{\textbackslash{}usepackage\{gb4e\}} to \texttt{header-includes}
in your YAML to enable glosses. Use \texttt{linguex} as an alternative
--- but not both, as they conflict.

Two-line gloss (\texttt{\textbackslash{}gll} \ldots{}
\texttt{\textbackslash{}glt}):

\begin{exe}
  \ex
  \gll Das Kind schläft. \\
       the child sleeps \\
  \glt `The child sleeps.'
\end{exe}

Three-line gloss with morpheme breaks (\texttt{\textbackslash{}glll}):

\begin{exe}
  \ex
  \glll un-believ-able \\
       \textsc{neg}-believe-\textsc{adj} \\
       un-believ-able \\
  \glt `unbelievable'
\end{exe}

Cross-reference examples with \texttt{\textbackslash{}Next},
\texttt{\textbackslash{}NNext}, or the \texttt{\textbackslash{}ref\{\}}
mechanism provided by the package.

\section{Masked Submission}\label{masked-submission}

Set \texttt{masked:\ true} in the YAML front matter to enable
double-blind mode. This suppresses all author and affiliation
information and replaces it with an automatic word count provided by the
bundled Lua filter. The \texttt{shorttitle}, \texttt{shortauthor}, and
\texttt{orcid} fields are also suppressed in masked mode.

Remember to remove or anonymise any self-identifying references in the
manuscript body before submitting. This includes references to your own
prior work that would reveal authorship, links to personal repositories,
and acknowledgements. Switch back to \texttt{masked:\ false} for the
final camera-ready version.

\section*{References}\label{references}
\addcontentsline{toc}{section}{References}

\printbibliography[heading=none]


\end{document}
